% John Dowland to Sir Robert Cecil (1595): study edition.
% Part I : full letter, original | modern English, the seven performance
%          passages highlighted, followed by the explanatory notes.
% Part II: one page per passage, original | modern English, with a list of
%          pronunciations and meanings at the bottom.
%
% Build (LuaLaTeX is required for the lua-ul highlighting), from documents/:
%   cd src && lualatex -output-directory=../out/doulande dowland-letter.tex
\documentclass[11pt,a4paper]{article}

\usepackage[top=18mm,bottom=20mm,left=17mm,right=17mm]{geometry}
\usepackage{fontspec}
\setmainfont{EBGaramond}[Extension=.otf, UprightFont=*-Regular, ItalicFont=*-Italic,
  BoldFont=*-SemiBold, BoldItalicFont=*-SemiBoldItalic]
\usepackage{microtype}
\usepackage[dvipsnames]{xcolor}
\usepackage{luacolor}
\usepackage{lua-ul}
\usepackage{paracol}
\usepackage{multicol}
\usepackage[hidelinks]{hyperref}

\setlength{\parindent}{0pt}
\setlength{\parskip}{0.55em}
\setlength{\columnsep}{7mm}
\setlength{\columnseprule}{0.3pt}
\def\columnseprulecolor{\color{black!25}}
\setlength{\emergencystretch}{2em}
\pagestyle{plain}
\frenchspacing

% ------------------------------------------------------------ passages
% Tag colour (dark, for labels) and highlight colour (light, behind text).
\definecolor{tag1}{HTML}{3F7D3A}  \colorlet{hl1}{tag1!22}
\definecolor{tag2}{HTML}{B8860B}  \colorlet{hl2}{tag2!25}
\definecolor{tag3}{HTML}{2C7FB8}  \colorlet{hl3}{tag3!20}
\definecolor{tag4}{HTML}{B03A5B}  \colorlet{hl4}{tag4!18}
\definecolor{tag5}{HTML}{6A4C9C}  \colorlet{hl5}{tag5!20}
\definecolor{tag6}{HTML}{3B4A5A}  \colorlet{hl6}{tag6!18}
\definecolor{tag7}{HTML}{1F8A86}  \colorlet{hl7}{tag7!22}
\colorlet{notecol}{black!55}

\newcommand{\ptitle}[1]{\ifcase#1\or
  O Sweet Woods, the Delight of Solitariness\or
  Fine Knacks for Ladies\or
  Clear or Cloudy, Sweet as April Showering\or
  Lasso vita mia, mi fa morire\or
  Come Again, Sweet Love doth Now Invite\or
  In Darkness Let Me Dwell\or
  Weep You No More, Sad Fountains\fi}

% Small coloured label that opens a passage.
\newcommand{\pbadge}[1]{\colorbox{tag#1}{\color{white}\bfseries\small\,#1\,}}
\newcommand{\ptag}[1]{\par\noindent\pbadge{#1}\enspace
  {\color{tag#1}\small\itshape\ptitle{#1}}\par\nobreak\vspace{-0.3em}}
% Highlight a stretch of original text in a passage colour.
\newcommand{\mk}[2]{\highLight[hl#1]{#2}}
% Word crossed out in the manuscript.
\newcommand{\strike}[1]{\strikeThrough{#1}}
% Note reference in the translation.
\newcommand{\fn}[1]{\textsuperscript{\color{notecol}#1}}
% Interlude inside passage 6.
\newcommand{\intermission}{\par\begin{center}\color{tag6}\small\itshape
  \textasciitilde\ [Marenzio song] \textasciitilde\end{center}\par}
\newcommand{\colhead}[1]{{\small\scshape\color{black!60}#1}\par\vspace{-0.2em}}

% Explanatory notes.
\newcommand{\note}[3]{\par\noindent{\bfseries #1\quad #2}\par\nobreak\vspace{-0.4em}#3\par}

% Glossary entry: word, pronunciation (may be empty), meaning.
\newcommand{\gl}[3]{\par\noindent\hangindent=1.2em\hangafter=1
  {\bfseries #1}\if\relax\detokenize{#2}\relax\else\ {\color{black!60}[#2]}\fi
  \ -- #3\par}

% John Dowland to Sir Robert Cecil, Nuremberg, 10 November 1595.
% Original text as transcribed at walterbitner.com (after Pinto); words the
% manuscript crosses out are shown with \strike. Modern English translation
% from john-dowland-letter-modern-english.docx.
%
% The letter is cut into pieces so each marked passage (\Oi ... \Ovii and
% \Ti ... \Tvii) can be used both in the full letter and on its own page.
% \OgX / \TgX are the unmarked stretches between passages.
% Passages in letter order: 1, 5, 2, 4, 3, 6, 7.

% ---------------------------------------------------------------- original

\newcommand{\OiA}{Ryght honorable: as I have bin most bounde unto your honor, so I most humblie desire your honor to pardon my boldnes, and mak my choic of your honor to let you understand my bounden duty and desire of Gods preservation of my most dear soveraigne Quene and contry: whom I besech God ever to bless and to confounde all their enemies what and whomsoever.}
\newcommand{\OiB}{Fiften yeares since I was in France Servant to Sir Henry Cobham who was Imbassador for the Quenes majestie, and lay in Parys, wher I fell aquainted with on Smith a preist, and on Morgan somtims of her majesties chapell, on Verstigan who brake out of Englande beinge aprehended, and on Moris a Welchman that was our porter, who is at Rome, these men thrust many idle toies into my hed of relygion, sainge that the papists was the truthe and ours in England all falce, and I being but yonge their faire words overecht me and I beleved with them.}

\newcommand{\OgA}{Within ij years after I cam into Englande wher I saw men of that faction condemde and executed, which I thought was great injustice taking relygion for the only cause, and when my best frends wold perswade me I wold not beleve them.}

\newcommand{\Ov}{Then in time passing on Mr.\ Johnson died, and I becam an humble sutor for his place (thincking my selfe most worthiest), wherin I found many goode and honorable frends that spake for me, but I saw that I was like to goe without it, and that any mygt have preferment but I. Wherby I began to sounde the cause, and gest that my relygion was my hinderance. Whearupon, my mynde beinge trobled, I desired to get beyond the seas,}

\newcommand{\OgB}{which I durst not attempt without lycence from som of the Privie Counsell, for fear of being taken and so my hav extreame punishment.}

\newcommand{\OiiA}{And accordinge as I desired ther cam a letter to me out of Germany from the Duke of Brunswicke, wherupon I spake to your honor and to my Lord of Essex, who willingly gav me both your hands (for which I wold be glad if ther wear any service in me that yor honors coulde comande).}
\newcommand{\OiiB}{When I cam to the Duke of Brunswicke he usde me kindly, and gave me a ritch chaine of golde, xxiij~£ in mony with velvet and saten and gold lace to make me aparell, with promyse that if I wolde serv him he wolde gev me as mutch as any prince in the worlde.}
\newcommand{\OiiC}{From thenc I went to the Landgrave of Hessen, (who gav me the gretest welcome that myght be for on of my qualyty) who sent a ringe into England to my wiff valued at xx~£ sterlinge, and gave me a great standing cupe with a cover gilt, full of dolers with many great offers for my service.}
\newcommand{\OiiD}{And from thence I had great desire to see Italy and cam to Venyce and from thence to Florence whear I plaid befor the Duke and got great favors.}

\newcommand{\Oiv}{And on evenynge I was walkinge upon the piazzo in Florence a jentle man tolde me that he espied a English preist and that his name was Skidmore and sonne and heir to Sir John Skidmore of the Courte, so I being intended to go unto Rome to study with a famous musicion named Luca Marenzio, stept to this Mr.\ Skidmor the preist and asked him if he was an Englysh man, and he told me yea: and whose sonne he was, and I telling him my name he was very glad to se me, so I tolde him I wold goe to Rome and desired his help for my saffty. For said I, “if they shold mystake ther, my fortune wear hard, for I hav bin thrust off of all good fortunes because I am a Catholicke at home.}

\newcommand{\Oiii}{For I hard that Her Majestie beinge spake to for me, saied, I was a man to serve any princ in the world, but I was an obstinat papist.” Whearunto he answered, “Mr.\ Dowlande, if it be not so, make her words true.” So in farther talke we spake of preists, and I told him that I did not thinck \strike{thinck} it true that any preists (as we said in England) wolde kill the Quen, or onc goe about to touche her finger, and said I whatsoever my Relygion be will neither medle nor mak with any thyng thear done, so that they do not any thing against the Quen. Whearunto he answered that I spake as a good Subiect to her majestie, but said he in Rome you shall Englyshmen your own Contrymen speak most hardly of her and wholy to overthrow her and all Englande, and those be the Jesuits, said he, who are of the Spanyshe faction: moreover said he we hav many Jares with them, and withall wished to God the Quen wear a Catholycke.}

\newcommand{\OgCa}{And said he, “to defend my Contry against the Spanyards I wold com into England and bear a pike on my sholders.”}
\newcommand{\OgCb}{Amonge our talk he tolde me that he had order to atache \strike{div} dyvers Inglish Jentlemen, and that he had bine iij years from Englande. So I brought him to his lodging dore, whear he told me that thear was ix prests com from Rome to go for Englande. He cam but the day before to Florenc, and I thinck they cam altogether.}
\newcommand{\OgCc}{He told me that he wold stay ther in the town and study in an abby caled Sancta Maria Nouella, and that he must kep in for on monethe, and that he wold writ letters of me to Rome, which I shold receve very shortly. But I heard not of him in a moneth after, and then there came two English friers to my lodging the on was an Englishman named Balye, a Yorksheirman.}
\newcommand{\OgCd}{The next day after my spech with Skydmore, I dyned with my Lord Gray and dyvers other jentlemen, whom I told of my speach with Skidmor, gevinge them warnynge. Wherupon my Lord Gray went to Siena, and the rest dispersed them selves.}
\newcommand{\OgCe}{More over I tolde my Lord Gray howsoever I was for relygion, if I did perceve anything in Rome that ether toucht her majestie of the state of England I wold gev nottice of it thoughe it wer the loss of my liff. Which he liked well and bade me kepe that secret.}
\newcommand{\OgCf}{This Frier Baylie before named delivered me a letter which I hav hear sent unto your honor, which letter I brake open befor Mr.\ \strike{bodly} Josias Bodly, and shewed what was written in it to him and divers other, after this, this frier Bayly tolde he had Receved letters From Rom to hasten me forward, and tolde me that my discontentment was knowne at Rom, and that I sholde hav a large pention of the Pope, and that His Holynes and all the Cardinales woulde make wonder full mutch of me.}
\newcommand{\OgCg}{Therupon I tolde him of my wiff and children how to get them to me, wherunto he tolde me that I shold have acquaintance with sutche as sholde bringe them over if she had any wilingnes, or elce they wolde lose ther lives. For ther cam those into Englande for sutch purposes. “For,” quothe he “Mr.\ Skidmor brought out of Englande at his last being ther xvij persons both men and women, for which the Bishope wepes when he sees him for Joy.”}

\newcommand{\OviA}{After my departure I caled to mynde our conference and got me by my self and wept hartely, to se my fortun so harde that I sholde becom sarvant to the gretest enemye of my prince, contry, wiff, children, and frends, for wante, and to make me like them selves.}
\newcommand{\OviB}{God He knoweth I never loved treason nor trechery nor never knew any, nor never heard any Mass in Englande, which I finde is great abuse of the peple for on my soule I understande it not, wherfor I have reformed my self to lyve acording to Her Majesties lawes as I was borne under Her Highnes, and that most humbly I doe crave pardon}

\newcommand{\OgDa}{, protestinge if ther wer any abylitie in me, I wold be most redy to make amends.}
\newcommand{\OgDb}{At Bolona I met with ij men the on named pierce an Iryshman, the other named Dracot, they are gon both to Rome, In Venice I hard an Italyan say, that he marveled that king Phillip had never a good frend in Inglande that with his dager wolde dispatch the Quens Majestie, but said he, God suffers her, in the end to gev her the greter overthrowe.}
\newcommand{\OgDc}{Right honorable, this hav I written that Her Majestie may knowe the vilany of thes most wicked prests and Jesuits, and to bewar of them. I thanck God I hav both forsaken them and their relygion which tendeth to nothing but distruction.}
\newcommand{\OgDd}{Thus I beseche God nyght and day to blesse and defend the Quens Majestie, and to confonde all her enemies and to preserve your honor and all the rest of her majesties most honorable Privie Counsell.}
\newcommand{\OgDe}{I thinck that Skidmore and the other preists ar all in England, for he staid not at Florenc as he said he wolde to me, and Frier Baylie tolde me that he was gon into France to study the lawe.}
\newcommand{\OgDf}{At Venyce and all along as I com in Germany men say that the Kinge of Spain is making gret preparation to com for England this next somer, wher}

\newcommand{\Ovii}{if it pleasde your honor to advise me by my poore wyff I wolde most wilingly lose my lyffe against them, most humbly besechinge your honor to pardon my ill writinge, and worse inditinge, and to thinke that I desir to serve my contry, and hope to hear of your good opinyon of me.}
\newcommand{\OviiDate}{From Nurnberge, this xth of November 1595}
\newcommand{\OviiClose}{Your honors most bounden}
\newcommand{\OSign}{for ever\\John Doulande}

% ------------------------------------------------------------- translation

\newcommand{\TiA}{Right Honourable Sir: Because I am deeply indebted to Your Honour, I most humbly ask you to forgive my presumption, and I choose Your Honour as the person to whom I may explain my solemn obligation and my desire that God preserve my most dear sovereign, the Queen, and my country.\fn{1} I pray that God will always bless them and confound every one of their enemies, whoever they may be.\fn{2}}
\newcommand{\TiB}{Fifteen years ago I was in France, in the service of Sir Henry Cobham, who was Her Majesty the Queen's ambassador and was residing in Paris.\fn{3} There I became acquainted with a priest named Smith; with Morgan, formerly a member of Her Majesty's Chapel; with Verstegan, who escaped from England after being apprehended; and with Morris, a Welshman who was our porter and who is now in Rome.\fn{4} These men filled my head with many foolish ideas about religion, saying that the papists possessed the truth and that our religion in England was entirely false. Because I was still young, their persuasive words deceived me, and I came to share their beliefs.\fn{5}}

\newcommand{\TgA}{Within two years after that I returned to England, where I saw men of that party condemned and executed. I believed this to be a great injustice, because I thought religion was the only reason for their punishment; and when my closest friends tried to persuade me otherwise, I would not believe them.\fn{6}}

\newcommand{\Tv}{As time passed, Mr Johnson died, and I humbly applied for his position, thinking myself the most worthy candidate.\fn{7} I found many good and honourable friends who spoke on my behalf, but I saw that I was likely to go without the position and that anyone might receive advancement except me. I therefore began to investigate the cause and guessed that my religion was the obstacle.\fn{8}

My mind being troubled by this, I wished to go abroad.}

\newcommand{\TgB}{I did not dare attempt it without permission from members of the Privy Council, for fear that I might be arrested and severely punished.\fn{9}}

\newcommand{\TiiA}{Just as I wished, a letter then came to me from Germany from the Duke of Brunswick. I therefore spoke to Your Honour and to my Lord of Essex, and you both willingly gave me your signatures and authorization.\fn{10} For this kindness I would be glad if there were any service within my power that Your Honours could command me to perform.}
\newcommand{\TiiB}{When I came to the Duke of Brunswick, he treated me kindly. He gave me a rich gold chain, twenty-three pounds in money, and velvet, satin and gold lace from which to make clothing. He also promised that, if I would serve him, he would give me as much as any prince in the world would give me.\fn{11}}
\newcommand{\TiiC}{From there I went to the Landgrave of Hesse, who gave me as great a welcome as could be given to a man of my profession and standing.\fn{12} He sent my wife in England a ring valued at twenty pounds sterling and gave me a large standing cup with a gilded cover, filled with thalers, together with many generous offers for my service.\fn{13}}
\newcommand{\TiiD}{From there I felt a great desire to see Italy. I came to Venice and then to Florence, where I played before the Duke and received great favour.\fn{14}}

\newcommand{\Tiv}{One evening, while I was walking in the piazza in Florence, a gentleman told me that he had spotted an English priest. He said that the priest's name was Skidmore and that he was the son and heir of Sir John Skidmore, a man of the royal court.\fn{15}

Because I intended to go to Rome to study with a famous musician named Luca Marenzio, I went up to Mr Skidmore, the priest, and asked whether he was an Englishman.\fn{16} He said that he was, and told me whose son he was. When I gave him my name, he was very glad to see me. I told him that I wished to go to Rome and asked for his help in securing my safety there.

I said, “If they should mistake me for something I am not when I arrive, my fortune would be hard indeed, because at home I have been excluded from every good opportunity on account of being considered a Catholic.\fn{17}}

\newcommand{\Tiii}{I have heard that, when Her Majesty was approached on my behalf, she said that I was a man capable of serving any prince in the world, but that I was an obstinate papist.”\fn{18}

He answered, “Mr Dowland, if that is not true, then make her words true.”\fn{19}

As our conversation continued, we spoke about priests. I told him that I did not believe it was true that any priests, as people claimed in England, would kill the Queen or even attempt to touch one of her fingers. I added, “Whatever my religion may be, I will neither involve myself nor take part in anything done there, provided that they do nothing against the Queen.”\fn{20}

He replied that I spoke like a loyal subject of Her Majesty. “But in Rome,” he said, “you will hear Englishmen, your own countrymen, speak most harshly of her and seek entirely to overthrow both her and all England. Those men are the Jesuits, who belong to the Spanish faction.” He added that his own group had many quarrels with them. At the same time, he wished to God that the Queen were a Catholic.}

\newcommand{\TgCa}{He then said, “To defend my country against the Spaniards, I would return to England and carry a pike on my shoulders.”\fn{21}}
\newcommand{\TgCb}{During our conversation, he told me that he had orders to apprehend various English gentlemen and that he had been away from England for three years.\fn{22} I accompanied him to the door of his lodging, where he told me that nine priests had come from Rome on their way to England. He himself had arrived in Florence only the previous day, and I think that they had all arrived together.}
\newcommand{\TgCc}{He told me that he would remain in the town and study in an abbey called Santa Maria Novella; that he must stay secluded there for one month; and that he would write letters about me, or on my behalf, to Rome, which I should receive very shortly.\fn{23} But I heard nothing from him for a month. Then two friars came to my lodging. One of them was an Englishman named Baylie, a Yorkshireman.\fn{24}}
\newcommand{\TgCd}{The day after my conversation with Skidmore, I dined with Lord Grey and several other gentlemen. I told them what Skidmore and I had discussed and warned them. Lord Grey consequently went to Siena, and the others dispersed.\fn{25}}
\newcommand{\TgCe}{I also told Lord Grey that, whatever my religion might be, if I discovered anything in Rome that affected Her Majesty or the English state, I would report it even if doing so cost me my life. He approved of this and instructed me to keep it secret.}
\newcommand{\TgCf}{The Friar Baylie mentioned above delivered a letter to me, which I have enclosed and sent to Your Honour. I broke the letter open in the presence of Mr Josias Bodley and showed him and several others what was written in it.\fn{26} After this, Friar Baylie told me that he had received letters from Rome urging me to proceed quickly. He said that my dissatisfaction was known in Rome, that I would receive a large pension from the Pope, and that His Holiness and all the cardinals would make a great deal of me.\fn{27}}
\newcommand{\TgCg}{I then asked him about my wife and children and how they could be brought to me. He told me that I would be introduced to people who could bring them over if my wife were willing; otherwise, those people would lose their lives, for men came into England for such purposes.\fn{28} “For,” he said, “when Mr Skidmore was last in England, he brought seventeen people out of the country, both men and women. The bishop weeps for joy whenever he sees him because of it.”\fn{29}}

\newcommand{\TviA}{After I left him, I recalled our conversation, withdrew by myself and wept bitterly. I grieved that my fortune should be so hard that poverty might force me to become the servant of the greatest enemy of my sovereign, my country, my wife, my children and my friends, and that they should try to make me like themselves.\fn{30}}
\newcommand{\TviB}{God knows that I have never loved treason or treachery, nor have I ever known of any. I have never attended Mass in England, which I now find to be a great deception of the people; upon my soul, I do not understand it.\fn{31} Therefore, I have reformed my conduct and resolved to live according to Her Majesty's laws, since I was born under Her Highness.\fn{32} For my past conduct I most humbly ask pardon,}

\newcommand{\TgDa}{declaring that, if I possess any ability to do so, I will be entirely ready to make amends.}
\newcommand{\TgDb}{At Bologna I met two men, one named Pierce, an Irishman, and the other named Dracot. They have both gone to Rome.\fn{33} In Venice I heard an Italian say that he was astonished King Philip had no good supporter in England who would dispatch Her Majesty the Queen with a dagger. He then said that God was allowing her to survive so that, in the end, He might give her an even greater overthrow.\fn{34}}
\newcommand{\TgDc}{Right Honourable Sir, I have written this so that Her Majesty may know the villainy of these most wicked priests and Jesuits and may be on her guard against them. I thank God that I have forsaken both them and their religion, which leads to nothing but destruction.\fn{35}}
\newcommand{\TgDd}{Thus I pray to God night and day to bless and defend Her Majesty the Queen, to confound all her enemies, and to preserve Your Honour and all the other members of Her Majesty's most honourable Privy Council.}
\newcommand{\TgDe}{I think that Skidmore and the other priests are now all in England, because he did not remain in Florence as he had told me he would. Friar Baylie told me that Skidmore had gone to France to study law, but I remain suspicious.\fn{36}}
\newcommand{\TgDf}{In Venice, and everywhere along my journey through Germany, people say that the King of Spain is making great preparations to come against England next summer.}

\newcommand{\Tvii}{If it pleased Your Honour to send me instructions through my poor wife, I would most willingly lose my life fighting against them.\fn{37}

I most humbly ask Your Honour to pardon my poor handwriting and my still poorer composition, and to believe that I desire to serve my country. I hope to hear that you hold a good opinion of me.\fn{38}}
\newcommand{\TviiDate}{From Nuremberg, this 10th day of November 1595}
\newcommand{\TviiClose}{Your Honour's most deeply obliged servant}
\newcommand{\TSign}{forever\\John Dowland}

% Introductory and closing texts, from john-dowland-letter-modern-english.docx.

\newcommand{\AboutText}{This version keeps close to Dowland’s meaning and sequence while replacing sixteenth-century spelling, grammar and idiom with current English. Where the original permits more than one reading, the translation chooses the most likely sense and records the uncertainty in the numbered notes. The notes also explain people, institutions and expressions that are easy to misread today.}

\newcommand{\DoingText}{The letter is simultaneously a confession, a defence, an intelligence report and an indirect request for renewed trust and patronage. Dowland constructs a careful story in which youthful persuasion led him toward Catholicism; his religious reputation blocked advancement in England; foreign princes valued him; Catholic agents then offered him money, honour and help for his family; and the discovery of anti-Elizabethan activity brought him back to open loyalty. His information may be sincere, but its presentation is strategic: it explains his past, distances him from treason and asks Cecil to regard him again as a useful English subject.}

\newcommand{\SourceText}{The translation follows David Pinto’s diplomatic and expanded transcription of the letter in the critical hypertext edition hosted by the University of Birmingham. The manuscript is held at Hatfield House in the Cecil Papers, volume 172, numbers 91-93.}

% Pronunciation and meaning lists for the seven passages (Part II).
% \gl{word as spelled}{say it as / stress in capitals}{meaning}
% Leave the middle argument empty when the spelling is the only difficulty
% and the meaning says it all.

\newcommand{\GlossI}{%
\gl{Ryght honorable}{right HON-ourable}{formal address to a Privy Councillor (Cecil)}
\gl{bin most bounde}{been most bound}{I am deeply indebted / obliged}
\gl{boldnes}{boldness}{presumption, forwardness}
\gl{mak my choic of your honor}{}{choose you as the one to tell}
\gl{bounden duty}{BOUN-den}{solemn moral obligation}
\gl{preservation}{preser-VAY-ceon}{protection, keeping safe}
\gl{soveraigne Quene}{SOV-rin queen}{Elizabeth I}
\gl{besech}{be-SAYCH}{beseech, beg}
\gl{confounde}{con-FOUND}{defeat, overthrow (not “confuse”)}
\gl{enemies}{ENEM-eyes}{enemies}
\gl{what and whomsoever}{}{whatever and whoever they are}
\gl{Fiften yeares since}{fifteen YAARS since}{fifteen years ago}
\gl{Sir Henry Cobham}{KOB-um}{English ambassador in Paris 1579–83}
\gl{Imbassador}{am-BASS-a-dor}{ambassador}
\gl{lay in Parys}{}{was residing in Paris}
\gl{fell aquainted with on Smith}{}{got to know one Smith (“on” = one)}
\gl{somtims of her majesties chapell}{}{formerly of the Chapel Royal}
\gl{Verstigan}{VER-ste-gan}{Richard Verstegan, Catholic writer and exile}
\gl{brake out \dots\ beinge aprehended}{}{escaped after being arrested}
\gl{Moris a Welchman}{MOR-ris, WELSH-man}{Morris, a Welshman}
\gl{idle toies}{idle toys}{foolish notions, fancies}
\gl{sainge}{SAY-ing}{saying}
\gl{papists}{PAY-pists}{Catholics (hostile Protestant term)}
\gl{falce}{false}{false}
\gl{yonge}{young}{young}
\gl{faire words}{}{persuasive, flattering talk}
\gl{overecht}{over-REACHT}{overreached: outwitted, deceived}
\gl{beleved with them}{}{came to share their beliefs}
}

\newcommand{\GlossII}{%
\gl{accordinge as I desired}{}{just as I wished}
\gl{ther cam}{there came}{there came}
\gl{Duke of Brunswicke}{BRUNZ-wick}{Heinrich Julius, court at Wolfenbüttel}
\gl{spake}{spake}{spoke}
\gl{my Lord of Essex}{ESS-ex}{Robert Devereux, Earl of Essex}
\gl{gav me both your hands}{}{gave me your signatures / authorisation}
\gl{service}{SAAR-vice}{service}
\gl{yor honors coulde comande}{}{that Your Honours could command}
\gl{usde me kindly}{yoozd}{treated me kindly}
\gl{ritch chaine of golde}{}{rich gold chain (mark of princely favour)}
\gl{xxiij £}{three-and-twenty pounds}{£23}
\gl{saten}{SAT-in}{satin}
\gl{aparell}{a-PA-rel}{clothing}
\gl{serv}{saarv}{serve}
\gl{gev \dots\ mutch}{give \dots\ much}{give \dots\ much}
\gl{worlde}{warld}{world}
\gl{From thenc}{from thence}{from there}
\gl{Landgrave of Hessen}{LAND-grave of HESS-en}{Moritz of Hesse-Kassel, himself a musician}
\gl{on of my qualyty}{oown of my QUAL-i-ty}{a man of my profession and rank}
\gl{wiff}{wife}{wife}
\gl{xx £ sterlinge}{twenty pounds STER-ling}{£20}
\gl{standing cupe with a cover gilt}{}{footed ceremonial cup with a gilded lid}
\gl{dolers}{DOL-ers}{thalers (TAH-lers), German silver coins}
\gl{Venyce}{VEN-iss}{Venice}
\gl{the Duke}{}{Ferdinando I de’ Medici (MED-i-chee), Grand Duke of Tuscany}
\gl{plaid befor}{played before}{performed for}
\gl{got great favors}{}{was received with great favour}
}

\newcommand{\GlossIII}{%
\gl{hard}{heard}{heard}
\gl{beinge spake to for me}{}{being petitioned on my behalf}
\gl{princ}{prince}{prince, ruler}
\gl{obstinat papist}{OB-sti-nate PAY-pist}{stubborn Catholic}
\gl{Whearunto}{where-un-TOO}{to which}
\gl{make her words true}{}{become what she called you}
\gl{farther talke}{}{further conversation}
\gl{thinck \strike{thinck}}{}{the repeated word is crossed out in the manuscript: say it once}
\gl{onc goe about to touche her finger}{oownce go about}{even try to lay a finger on her}
\gl{medle nor mak}{meddle nor make}{neither interfere nor take part}
\gl{so that}{}{provided that}
\gl{Subiect}{SUB-ject}{subject (i written for j)}
\gl{you shall Englyshmen}{}{read “you shall [hear] Englishmen”: Dowland drops a word}
\gl{most hardly}{}{most harshly}
\gl{wholy to overthrow}{HOLE-ly}{seek entirely to overthrow}
\gl{Jesuits}{JEZ-yoo-its}{Society of Jesus, here the pro-Spanish exiles}
\gl{Spanyshe faction}{SPAN-ish FACT-ceon}{party aligned with Philip II of Spain}
\gl{Jares}{jars}{quarrels}
\gl{withall}{with-ALL}{at the same time}
\gl{Catholycke}{CATH-o-lic}{Catholic}
}

\newcommand{\GlossIV}{%
\gl{on evenynge}{one EVE-ning}{one evening}
\gl{piazzo}{pee-AT-sa}{piazza, town square}
\gl{jentle man}{gentleman}{gentleman}
\gl{espied}{es-PIDE}{caught sight of}
\gl{Skidmore}{SKID-more}{John Scudamore (SKUD-a-more), English Catholic priest}
\gl{of the Courte}{}{of the royal court}
\gl{being intended}{}{intending}
\gl{musicion}{mew-ZISH-un}{musician}
\gl{Luca Marenzio}{LOO-ka ma-REN-tsee-o}{leading Italian madrigal composer, in Rome}
\gl{name / named}{naaym}{name}
\gl{stept to}{}{stepped up to}
\gl{Englysh man}{}{Englishman}
\gl{yea}{yay}{yes}
\gl{se}{see}{see}
\gl{saffty}{SAFE-ty}{safety}
\gl{if they shold mystake ther}{}{if they should misjudge me there (in Rome)}
\gl{wear hard}{were hard}{would be hard}
\gl{thrust off of all good fortunes}{aal good FART-ns}{shut out of every good opportunity}
\gl{a Catholicke at home}{}{(taken for) a Catholic in England}
}

\newcommand{\GlossV}{%
\gl{in time passing on}{}{as time went by}
\gl{on}{oown}{on}
\gl{Mr.\ Johnson}{}{John Johnson, royal lutenist, died 1594}
\gl{becam an humble sutor}{SOO-ter}{humbly applied, petitioned}
\gl{his place}{his plaayce}{his post in the royal music}
\gl{most worthiest}{most WARTH-yest}{the most worthy (double superlative)}
\gl{wherin}{where-IN}{in which (matter)}
\gl{spake for me}{}{spoke on my behalf}
\gl{like to goe without it}{}{likely to miss out on it}
\gl{mygt have preferment}{might have pre-FER-ment}{might receive appointment, advancement}
\gl{Wherby}{}{whereby, so}
\gl{sounde the cause}{}{probe, investigate the reason}
\gl{gest}{guessed}{guessed}
\gl{hinderance}{HIN-drance}{obstacle}
\gl{Whearupon}{where-up-ON}{whereupon, and so}
\gl{mynde beinge trobled}{}{my mind being troubled}
\gl{beyond the seas}{}{abroad, to the Continent}
}

\newcommand{\GlossVI}{%
\gl{caled to mynde}{called to mind}{recalled}
\gl{conference}{}{conversation}
\gl{got me by my self}{}{went off alone}
\gl{wept hartely}{HEART-ily}{wept bitterly, from the heart}
\gl{se my fortun so harde}{see my FART-n so hard}{to see my fortune so hard}
\gl{sarvant}{SAAR-vant}{servant}
\gl{gretest enemye}{}{greatest enemy (the Pope)}
\gl{my prince}{}{my sovereign: the Queen (“prince” = any ruler)}
\gl{for wante}{}{out of poverty, for lack of a living}
\gl{to make me like them selves}{}{that they should make me one of them}
\gl{[Marenzio song]}{}{interlude between the two halves}
\gl{God He knoweth}{}{God knows}
\gl{treason}{TRAY-son}{treason}
\gl{trechery}{TRECH-er-y}{treachery}
\gl{nor never knew any}{}{nor ever knew of any (double negative)}
\gl{Mass}{}{the Catholic Latin service}
\gl{great abuse of the peple}{a-BUSE}{a great deception of the people}
\gl{on my soule}{}{upon my soul}
\gl{wherfor}{where-FOR}{therefore}
\gl{reformed my self}{}{changed my ways, resolved to conform}
\gl{as I was borne under Her Highnes}{}{since I was born her subject}
\gl{crave pardon}{}{beg forgiveness}
}

\newcommand{\GlossVII}{%
\gl{pleasde}{pleased}{pleased}
\gl{advise me by my poore wyff}{}{send me instructions through my wife}
\gl{wilingly}{}{willingly}
\gl{lose my lyffe against them}{}{die fighting the Spanish}
\gl{besechinge}{be-SAY-ching}{beseeching, begging}
\gl{ill writinge}{}{bad handwriting}
\gl{worse inditinge}{in-DITE-ing}{even worse composition, wording}
\gl{desir}{de-ZIRE}{desire}
\gl{contry}{COUN-try}{country}
\gl{opinyon}{o-PIN-yun}{opinion}
\gl{Nurnberge}{NURN-berg}{Nuremberg, where the letter was written}
\gl{xth}{tenth}{10th (November 1595)}
\gl{Your honors most bounden}{}{Your Honour’s most indebted (servant)}
}


\begin{document}

% ================================================================= PART I
\begin{center}
  {\LARGE John Dowland to Sir Robert Cecil}\par\vspace{0.3em}
  {\large Nuremberg, 10 November 1595}\par\vspace{0.2em}
  {\small\color{black!60}Cecil Papers 172/91--93 \textperiodcentered\
  original text with a modern English translation, notes and the seven
  marked passages}
\end{center}

\vspace{0.4em}
{\small
\textbf{The seven passages}, in performance order. In the letter they occur
in the order 1, 5, 2, 4, 3, 6, 7; each is highlighted in its colour on the
original side.\par\vspace{0.2em}
\begin{multicols}{2}
\raggedright
\pbadge{1}\enspace\ptitle{1}\par
\pbadge{2}\enspace\ptitle{2}\par
\pbadge{3}\enspace\ptitle{3}\par
\pbadge{4}\enspace\ptitle{4}\par
\pbadge{5}\enspace\ptitle{5}\par
\pbadge{6}\enspace\ptitle{6}\par
\pbadge{7}\enspace\ptitle{7}\par
\end{multicols}
\vspace{-0.4em}
\textbf{About this translation.} \AboutText\par
Words struck through (\strike{thinck}) are crossed out in the manuscript.
Superscript numbers refer to the explanatory notes after the letter.\par}

\vspace{0.3em}
\section*{I.\enspace The letter}

\begin{paracol}{2}
\colhead{Original (1595)}
\switchcolumn
\colhead{Modern English}
\switchcolumn*

\ptag{1}\mk{1}{\OiA}
\switchcolumn
\ptag{1}\TiA
\switchcolumn*
\mk{1}{\OiB}
\switchcolumn
\TiB
\switchcolumn*
\OgA
\switchcolumn
\TgA
\switchcolumn*

\ptag{5}\mk{5}{\Ov}
\switchcolumn
\ptag{5}\Tv
\switchcolumn*
\OgB
\switchcolumn
\TgB
\switchcolumn*

\ptag{2}\mk{2}{\OiiA}
\switchcolumn
\ptag{2}\TiiA
\switchcolumn*
\mk{2}{\OiiB}
\switchcolumn
\TiiB
\switchcolumn*
\mk{2}{\OiiC}
\switchcolumn
\TiiC
\switchcolumn*
\mk{2}{\OiiD}
\switchcolumn
\TiiD
\switchcolumn*

\ptag{4}\mk{4}{\Oiv}
\switchcolumn
\ptag{4}\Tiv
\switchcolumn*

\ptag{3}\mk{3}{\Oiii}
\switchcolumn
\ptag{3}\Tiii
\switchcolumn*

\OgCa
\switchcolumn
\TgCa
\switchcolumn*
\OgCb
\switchcolumn
\TgCb
\switchcolumn*
\OgCc
\switchcolumn
\TgCc
\switchcolumn*
\OgCd
\switchcolumn
\TgCd
\switchcolumn*
\OgCe
\switchcolumn
\TgCe
\switchcolumn*
\OgCf
\switchcolumn
\TgCf
\switchcolumn*
\OgCg
\switchcolumn
\TgCg
\switchcolumn*

\ptag{6}\mk{6}{\OviA}
\switchcolumn
\ptag{6}\TviA
\switchcolumn*
\intermission
\switchcolumn
\intermission
\switchcolumn*
\mk{6}{\OviB}\OgDa
\switchcolumn
\TviB\ \TgDa
\switchcolumn*

\OgDb
\switchcolumn
\TgDb
\switchcolumn*
\OgDc
\switchcolumn
\TgDc
\switchcolumn*
\OgDd
\switchcolumn
\TgDd
\switchcolumn*
\OgDe
\switchcolumn
\TgDe
\switchcolumn*
\OgDf
\switchcolumn
\TgDf
\switchcolumn*

\ptag{7}\mk{7}{\Ovii}
\switchcolumn
\ptag{7}\Tvii
\switchcolumn*
{\itshape \OviiDate}\par
\mk{7}{\OviiClose}\\
\OSign
\switchcolumn
{\itshape \TviiDate}\par
\TviiClose\\
\TSign
\end{paracol}

% ------------------------------------------------------------ notes
\section*{Explanatory notes}
{\small
\begin{multicols}{2}
\setlength{\parskip}{0.35em}
% Explanatory notes 1-38, from john-dowland-letter-modern-english.docx.

\note{1}{Right Honourable and bounden duty}{Dowland addresses Robert Cecil, then a Privy Councillor. ‘Bounden’ means morally and personally obliged or indebted, not physically restrained. The opening combines conventional humility with a practical appeal to Cecil as patron and political protector.}

\note{2}{Preserve and confound}{To ask for God’s ‘preservation’ is to ask that God protect and sustain the Queen and country. ‘Confound’ here means frustrate, defeat or overthrow, rather than merely confuse. The pronoun ‘them’ refers jointly to Queen and country.}

\note{3}{Sir Henry Cobham and lay in Paris}{Cobham was the English ambassador in Paris from 1579 to 1583. Dowland’s being his ‘servant’ means that he belonged to the ambassador’s household or service; it need not mean that he performed ordinary domestic labour. To ‘lie’ in a city could mean to reside or be stationed there.}

\note{4}{The Paris acquaintances}{Nicholas Morgan had formerly been a gentleman of the Chapel Royal, an institution of royal musicians and clergy rather than simply a physical chapel. Richard Verstegan was a Catholic writer and exile. Morris was the household porter or doorkeeper; ‘at Rome’ was a normal idiom for residing in Rome.}

\note{5}{Idle toys and overreached}{‘Idle toys’ means foolish ideas, fantasies or trifles; ‘toys’ could mean notions, not only children’s objects. Dowland’s unusual spelling ‘overecht’ is best understood as ‘overreached’: they outwitted, misled or deceived him. ‘I believed with them’ means that he adopted or shared their beliefs.}

\note{6}{Within two years after}{The phrase means that within two years after his Paris experience Dowland came back to England, not that the executions occurred within two years after his arrival. ‘That faction’ means the Catholic party or group. Dowland says that at the time he believed religion alone caused the punishments; the letter later presents this belief as youthful error.}

\note{7}{Mr Johnson’s place}{John Johnson, one of the leading English lutenists, served in the royal musical establishment and died in 1594. Dowland became a ‘suitor’ for his place: an applicant or petitioner. His aside that he thought himself ‘most worthiest’ reveals both confidence in his ability and resentment at the outcome.}

\note{8}{Preferment and sounding the cause}{‘Preferment’ means appointment, promotion or advancement through patronage. To ‘sound’ the cause means to probe or investigate it. Dowland says he guessed that his reputation as a Catholic prevented his appointment; he does not prove that this was the official reason.}

\note{9}{Beyond the seas and the Privy Council}{‘Beyond the seas’ means abroad, especially continental Europe. ‘Durst’ is the old past tense of ‘dare’. The Privy Council was the Queen’s principal executive and advisory body; ‘privy’ means private or confidential. Unlicensed foreign travel could attract suspicion, particularly amid fears of Catholic conspiracy.}

\note{10}{Both your hands}{The phrase most likely means that Cecil and the Earl of Essex supplied their signatures, warrants or written approval. It does not mean that they physically gave Dowland both hands. His parenthesis offers future service in repayment.}

\note{11}{The Duke of Brunswick’s gifts}{The duke was Heinrich Julius of Brunswick-Luneburg, whose court was at Wolfenbuttel. The chain, money and luxury textiles were marks of princely favour and potential employment. ‘Apparel’ means clothing; historically the word is related to equipping or preparing a person.}

\note{12}{One of my quality}{Here ‘quality’ means rank, profession, accomplishment or social standing. Dowland is saying that the welcome was exceptionally great for a musician of his position, not commenting on the abstract quality of his music.}

\note{13}{The Landgrave’s gifts}{Moritz, Landgrave of Hesse-Kassel, was himself musically accomplished. A ‘standing cup’ was a large ceremonial drinking vessel with a foot and often a lid. The ‘dollars’ were thalers. The gifts to Dowland and his wife demonstrate the status he enjoyed abroad.}

\note{14}{The Duke in Florence}{This was Ferdinando I de’ Medici, Grand Duke of Tuscany. ‘Got great favours’ means that Dowland received a favourable reception and patronage, not necessarily one particular formal reward.}

\note{15}{Skidmore or Scudamore}{The priest was John Scudamore, whose name Dowland spells in several ways. He was an English Catholic priest, not an Anglican clergyman. The unnamed gentleman in the piazza merely pointed him out; the gentleman plays no further part in the story.}

\note{16}{Luca Marenzio}{Marenzio was one of the leading Italian madrigal composers of the period and was associated with Rome. Dowland says that musical study was the purpose of his intended journey. The letter does not say that Skidmore would travel with him, and Dowland apparently did not reach Rome on this journey.}

\note{17}{Mistaken in Rome}{Dowland probably fears being mistaken for a Protestant agent, spy or religious enemy in Catholic Rome. The wording is deliberately general. He presents himself as caught between two suspicions: treated as Catholic in England but potentially mistrusted in Rome.}

\note{18}{Obstinate papist}{‘Papist’ was a hostile Protestant term for a Roman Catholic. The Queen’s alleged statement is hearsay reported by Dowland, not a surviving direct quotation from Elizabeth. It praises his musical usefulness while identifying his religion as a political obstacle.}

\note{19}{Make her words true}{Skidmore’s reply means: if you are not yet an obstinate Catholic, become one and thereby make the Queen’s description accurate. It is an invitation to commit himself more fully to Catholicism.}

\note{20}{Touch her finger and meddle nor make}{To say that priests would not even ‘touch her finger’ intensifies Dowland’s disbelief that they would harm Elizabeth. ‘Neither meddle nor make with’ means neither become involved nor take part. His qualification is political: he will stay uninvolved so long as nothing is done against the Queen.}

\note{21}{Spanish faction jars and the pike}{The ‘Spanish faction’ means Catholics aligned politically with Philip II of Spain. ‘Jars’ are quarrels or conflicts. To ‘bear a pike’ means to serve as an ordinary infantryman. Skidmore presents himself as Catholic but English and anti-Spanish, illustrating divisions among exiled English Catholics.}

\note{22}{Order to attach}{In sixteenth-century usage ‘attach’ commonly meant arrest, seize or take into legal custody. That is the likeliest sense, although Dowland gives too little context to establish what authority Skidmore claimed or which gentlemen were meant.}

\note{23}{Keep in and write letters of me}{To ‘keep in’ probably means remain indoors, secluded or confined for a month. ‘Write letters of me’ can mean write about Dowland or write recommendations on his behalf. Santa Maria Novella was a major Dominican religious house in Florence.}

\note{24}{Friar Baylie}{A friar is a member of a Catholic mendicant religious order; the closest German terms are ‘Ordensbruder’ or, more specifically, ‘Bettelmönch’. It does not simply mean any Catholic priest, but Baylie is certainly presented as a Catholic religious.}

\note{25}{Warning Lord Grey}{Dowland says he immediately warned English gentlemen about Skidmore’s information, whereupon they separated. This episode helps him present himself to Cecil as already acting loyally before he wrote the letter.}

\note{26}{Josias Bodley as witness}{Dowland stresses that he opened Baylie’s letter before Josias Bodley and other witnesses. He also enclosed the letter for Cecil. These details provide corroboration and make the account look like an intelligence report rather than secret Catholic correspondence.}

\note{27}{The papal offer}{Baylie’s promises address precisely Dowland’s grievances: lack of employment, insufficient money and exclusion from honour. Whether a pension had truly been authorized is unknown; within Dowland’s narrative, the offer represents the height of the temptation.}

\note{28}{Who might lose their lives}{The pronoun is ambiguous. It most probably refers to the agents who entered England to escort Catholics out and who risked execution if caught. It could grammatically be taken to refer to Dowland’s family, but the following sentence about agents entering England favours the first reading.}

\note{29}{Seventeen people and the bishop}{This is Baylie’s claim, repeated by Dowland. The bishop is not identified. The story praises Skidmore for helping English Catholics escape, but Dowland now offers it to Cecil as evidence of clandestine activity.}

\note{30}{For want}{Here ‘want’ means poverty, need or lack of livelihood. Dowland says that professional exclusion might drive him into papal service. The passage is also carefully self-defensive: material desperation, rather than disloyal intention, has brought him close to England’s enemies.}

\note{31}{Treason treachery and Mass}{‘Treason’ is betrayal of the sovereign or state and had a specific legal force; ‘treachery’ is the broader idea of deceitful betrayal. The grammar most naturally makes Mass the ‘great abuse of the people’, perhaps because its Latin ritual was unintelligible to ordinary worshippers and, Dowland says, even to him. The sentence is polemical and may also be a calculated profession of Protestant sympathy.}

\note{32}{Reformed myself}{In context this certainly promises obedience to Elizabeth’s religious laws. Whether it proves an inward theological conversion is less certain. Dowland may be promising outward conformity and lawful conduct rather than describing a complete change of private belief.}

\note{33}{Pierce and Dracot}{Dowland gives no accusation or further story about these men. Their names may simply be intelligence leads for Cecil’s network. Their mention is relevant because both were travelling to Rome, but their precise role remains unknown.}

\note{34}{King Philip}{This is Philip II of Spain, Catholic king and former husband of Mary I of England. Spain and England were at war; the Spanish Armada had attempted an invasion in 1588, and further Spanish operations remained a serious fear. ‘Dispatch’ here means kill.}

\note{35}{Forsaking them and their religion}{This sounds like an explicit renunciation of Catholicism. Nevertheless, the letter as a whole is a plea to a powerful Protestant statesman, and its sincerity cannot be measured directly. The safer conclusion is that Dowland publicly rejects the Catholic network described here and promises conformity and loyalty to Elizabeth.}

\note{36}{Skidmore’s whereabouts}{Dowland infers that Skidmore and the priests may be in England because Skidmore did not remain in Florence as promised. Baylie’s statement that Skidmore went to France to study law is reported but apparently not trusted. Dowland is offering suspicion, not confirmed knowledge.}

\note{37}{Instructions through my wife}{Dowland asks Cecil to communicate with him through Dowland’s wife in England. His offer to ‘lose my life against them’ means to fight, and if necessary die, against a Spanish attack. It is chiefly a pledge of loyalty, not a clearly formulated request to enlist in a particular military unit.}

\note{38}{Ill writing and worse inditing}{‘Writing’ refers to handwriting; ‘inditing’ means composing or wording the letter. The conventional apology closes a broader request: Cecil should pardon Dowland’s presumption, accept his account, restore his reputation and consider him available for English service.}

\end{multicols}

\subsection*{What the letter is doing}
\DoingText

\subsection*{Textual source}
\SourceText\ Critical edition and transcription:
\url{https://philological.cal.bham.ac.uk/dowland/text.html}; commentary notes:
\url{https://philological.cal.bham.ac.uk/dowland/notes.html}. Original
spelling as given at
\url{https://walterbitner.com/john-dowland-to-sir-robert-cecil-november-10-1595/}.
\par}

% ================================================================= PART II
% One page per passage.
\newcommand{\passagepage}[4]{%
  \begin{paracol}{2}
  \colhead{Original (1595)}
  #2
  \switchcolumn
  \colhead{Modern English}
  #3
  \end{paracol}
  % Passage caption below the excerpt.
  \par\vspace{0.4em}
  \noindent{\Large\colorbox{tag#1}{\color{white}\bfseries\,#1\,}\enspace
    \color{tag#1}\ptitle{#1}}\par\vspace{0.1em}
  \noindent{\color{tag#1}\rule{\linewidth}{1.2pt}}\par
  \vfill
  \noindent{\color{tag#1}\rule{\linewidth}{0.6pt}}\par
  {\small\colhead{Pronunciation and meanings}
  \begin{multicols}{2}
  \setlength{\parskip}{0.15em}\raggedright
  #4
  \end{multicols}}
}

\clearpage
\section*{II.\enspace The seven passages}
Each passage on its own page: original on the left, modern English on the
right, and at the bottom a list of words to say or understand. Old spellings
are read as the modern word given in brackets; capitals mark the stressed
syllable. Several brackets give a period sound rather than the modern
one (e.g.\ \emph{saarv} for serve, \emph{warld} for world, \emph{oown} for
one, \emph{fart-n} for fortune).

\passagepage{1}{\OiA\ \OiB}{\TiA\par\TiB}{\GlossI}
\clearpage\passagepage{2}{\OiiA\ \OiiB\ \OiiC\ \OiiD}{\TiiA\par\TiiB\par\TiiC\par\TiiD}{\GlossII}
\clearpage\passagepage{3}{\Oiii}{\Tiii}{\GlossIII}
\clearpage\passagepage{4}{\Oiv}{\Tiv}{\GlossIV}
\clearpage\passagepage{5}{\Ov}{\Tv}{\GlossV}
\clearpage\passagepage{6}{\OviA\intermission\OviB}{\TviA\intermission\TviB}{\GlossVI}
\clearpage\passagepage{7}{\Ovii\par\OviiClose}{\Tvii\par\TviiClose}{\GlossVII}

\end{document}
